\chapter{落地执行、真实性证据与实施计划}

\section{当前完成度分层}

智舆当前成果采用已完成、开发中、待验证三层口径表达。该口径用于清晰展示项目成熟度。

\begin{table}[htbp]
	\t\centering
	\t\caption{智舆当前完成度分层}
	\t\begin{tabularx}{\textwidth}{|L{2.2cm}L{4.1cm}YL{3.3cm}|}
		\t\t\toprule
		\t\t\tblhead{状态} & \tblhead{主要内容}                & \tblhead{当前说明}                                     & \tblhead{证据形式}       \\
		\t\t\midrule
		\t\t\rowcolors{2}{tablealt}{white}
		\t\t已完成          & `paper3` 商业计划书主工程、章节骨架、核心内容框架 & 已形成可持续迭代的 LaTeX 文档结构，摘要、背景、产品、技术、市场、商业等主体章节已具备正文基础 & PDF 编译结果、章节截图、仓库目录截图 \\
		\t\t已完成          & 项目概念原型与业务闭环表达                 & 已明确“采集-分析-预警-推演-反馈”的业务闭环，具备产品展示和方案交流基础             & 原型截图、流程图、演示录屏        \\
		\t\t开发中          & 多场景试点素材与证据归档                  & 需要进一步整理界面截图、测试记录、用户反馈与沟通纪要，形成可核验附录包                & 测试报告、沟通纪要、复盘材料       \\
		\t\t开发中          & 面向试点的实施方案细化                   & 已形成试点路径与指标框架，但仍需结合具体单位数据源和流程边界做定制调整                & 试点方案书、需求确认表、实施计划表    \\
		\t\t待验证          & 本地正式合作材料与盖章文件                 & 合同、合作意向、授权文件、正式采购材料等需按真实进展补充，不在正文中提前写成已签约事实        & 合同扫描件、意向函、授权书        \\
		\t\t待验证          & 经营性指标与长期续费数据                  & 客户续费率、真实收入流水、跨区域复制效果等需在试点落地后逐步沉淀                   & 财务流水、客户回访记录、续费统计表    \\
		\t\t\bottomrule
		\t\end{tabularx}
\end{table}

\section{试点落地方案}

首轮试点优先围绕信阳本地高频、可验证、易形成成效指标的治理场景展开，采用小切口切入、短周期复盘、低门槛复制的原则。项目当前更适合从轻量化场景做样板，首轮优先选择景区文旅服务、校园管理、基层民生反馈和重点企业品牌舆情四类对象开展验证。

\begin{table}[htbp]
	\t\centering
	\t\caption{首轮试点场景}
	\t\begin{tabularx}{\textwidth}{|L{2.6cm}L{3.1cm}YL{2.8cm}|}
		\t\t\toprule
		\t\t\tblhead{试点场景} & \tblhead{典型对象} & \tblhead{核心价值}                    & \tblhead{优先输出}  \\
		\t\t\midrule
		\t\t\rowcolors{2}{tablealt}{white}
		\t\t文旅景区场景         & 景区运营方、文旅主管单位   & 监测投诉热点、节假日舆情波动、服务体验问题，缩短热点发现和回应链路 & 热点预警看板、节假日专题报告  \\
		\t\t校园治理场景         & 高校宣传或学工部门      & 监测校园热点事件、社媒情绪变化和突发话题扩散风险          & 校园周报、事件快报、应对建议单 \\
		\t\t基层民生场景         & 街道、热线协同单位      & 识别高频投诉主题、突发民生问题和跨渠道重复诉求           & 主题聚类、处置优先级清单    \\
		\t\t企业品牌场景         & 本地企业、公服机构      & 监测品牌声量、客户评价和舆论风险，辅助企业公关与服务优化      & 舆情日报、风险标签、传播趋势图 \\
		\t\t\bottomrule
		\t\end{tabularx}
\end{table}

\textbf{首个企业客户试点进展：}和鸣科技(武汉)有限公司已与\textbf{郑州浩龙机电设备有限公司}完成业务沟通，双方就企业舆情监测与品牌风险管理需求达成初步合作共识。浩龙机电作为制造业企业，在企业管理、品牌维护和客户反馈监测方面具有典型需求，拟采用本地部署模式开展首轮试点，当前正在推进需求细化与合同签署工作。该试点将为同类企业客户提供可复制的实施样板。

在试点策略上，项目不追求一开始就对接复杂的全量政务系统，而是先完成可见成效的小闭环。先把数据来源、预警阈值、处置流程和复盘机制跑通，再逐步增加接口集成、专题分析和多 Agent 推演模块，降低首批客户的理解和采购门槛。

\section{实施路径与里程碑}

为保证项目从展示型原型走向可交付方案，实施路径按照调研、接入、联调、培训、上线、评估六个阶段推进，并为每个阶段设置明确输出物。

\begin{table}[htbp]
	\t\centering
	\t\caption{智舆试点实施路径}
	\t\begin{tabularx}{\textwidth}{|L{2cm}L{2.1cm}YL{3cm}|}
		\t\t\toprule
		\t\t\tblhead{阶段} & \tblhead{周期} & \tblhead{关键任务}                  & \tblhead{阶段输出物}   \\
		\t\t\midrule
		\t\t\rowcolors{2}{tablealt}{white}
		\t\t需求调研         & 第 1 周        & 梳理使用场景、目标用户、数据来源、事件分级规则与响应链路    & 需求访谈表、场景清单、试点边界说明 \\
		\t\t数据接入         & 第 2--3 周     & 接入样例数据或公开数据，完成字段映射、清洗规则与主题标签初始化 & 数据接入清单、字段映射表、样例库  \\
		\t\t模型联调         & 第 3--4 周     & 调整分类标签、情绪识别阈值、告警策略与报告模板         & 调优记录、测试样例、报告模板    \\
		\t\t业务试运行        & 第 5--6 周     & 在限定范围内输出日报、周报、快报和预警提醒，验证闭环效率    & 试运行日报、问题台账、迭代清单   \\
		\t\t培训上线         & 第 6 周        & 对使用人员进行培训，明确权限、流程、人工复核和反馈机制     & 培训材料、操作手册、权限表     \\
		\t\t效果评估         & 第 7--8 周     & 统计响应时间、漏报情况、人工工时节约与满意度变化，形成复盘结论 & 评估报告、复盘纪要、下一阶段建议  \\
		\t\t\bottomrule
		\t\end{tabularx}
\end{table}

\section{试点成效指标设计}

试点成效不宜只用系统能否运行来判断，而应围绕治理效率、识别质量、协同效果和管理满意度四个维度进行评价。

\begin{table}[htbp]
	\t\centering
	\t\caption{试点成效指标}
	\t\begin{tabularx}{\textwidth}{|L{3cm}L{3cm}Y|}
		\t\t\toprule
		\t\t\tblhead{指标维度} & \tblhead{观察指标}         & \tblhead{说明}            \\
		\t\t\midrule
		\t\t\rowcolors{2}{tablealt}{white}
		\t\t预警效率           & 从舆情出现到系统预警的平均时长        & 用于评估系统是否真正缩短发现问题的时间     \\
		\t\t研判质量           & 事件分类准确率、情绪识别可用率、误报漏报情况 & 体现模型输出是否具备辅助决策价值        \\
		\t\t协同效率           & 从发现问题到完成责任分派的平均时长      & 体现系统是否推动跨岗位协同，而非停留在看板展示 \\
		\t\t人工节约           & 每周减少的人工检索、汇总、写报工时      & 便于量化 AI+ 带来的效率提升        \\
		\t\t管理满意度          & 试点单位使用人员评分与复盘意见        & 体现产品是否被一线用户真正接受         \\
		\t\t复制潜力           & 试点后可复制模块数量与标准化程度       & 判断项目是否具备从单点试点走向批量复制的能力  \\
		\t\t\bottomrule
		\t\end{tabularx}
\end{table}

\section{从原型到产品化的推进路线}

从原型产品走向可持续产品，需要完成三个层级的跃迁。从可展示走向可使用，需要把关键页面、数据链路和报告模板打磨到真实用户能理解、能操作。从可使用走向可交付，需要补齐权限管理、实施文档、培训流程和复盘机制。从可交付走向可复制，需要将首批试点中反复出现的共性模块抽象为标准包，并沉淀区域适配模板、行业标签库和服务手册。

\section{真实性证据目录}

附录按界面材料、调研材料、测试材料、试点材料、经营材料五个大类归档，统一采用附件编号方式管理，例如 E-01 表示界面类证据、R-01 表示调研类证据、T-01 表示测试类证据。该编号方式便于快速定位关键证明材料。
